\documentclass[12pt,a4paper]{article}

% ------------------------------------------------
% PACKAGES
% ------------------------------------------------
\usepackage[margin=1in]{geometry}
\usepackage{setspace}
\usepackage{titlesec}
\usepackage{enumitem}
\usepackage{fancyhdr}
\usepackage{hyperref}
\usepackage{booktabs}
\usepackage{longtable}
\usepackage{lastpage}
\usepackage{array}

\setstretch{1.15}

% ------------------------------------------------
% HEADER / FOOTER
% ------------------------------------------------
\pagestyle{fancy}
\fancyhf{}
\lhead{IBDP Psychology (2025)}
\rhead{Paper 1 Framework}
\rfoot{Page \thepage\ of \pageref{LastPage}}

\hypersetup{
    colorlinks=true,
    linkcolor=black
}

% ------------------------------------------------
% DOCUMENT
% ------------------------------------------------
\begin{document}

\begin{center}
    {\Large \textbf{IBDP Psychology (First Assessment 2025)}}\\[6pt]
    {\large Paper 1 – Assessment Framework (Internal Guide)}\\[10pt]
    \rule{\textwidth}{0.4pt}
\end{center}

\vspace{0.5cm}

\section*{Purpose}
This document provides a structured overview of Paper 1 under the 2025 assessment model. 
It is designed for instructional planning and internal academic use. 
The content reflects the structure and expectations of the assessment while using original wording.

\newpage

% ------------------------------------------------
\section{Paper 1 Overview}

\begin{itemize}
    \item \textbf{Duration:} 1 hour 30 minutes
    \item \textbf{Total Marks:} 35
    \item \textbf{Weighting:}
    \begin{itemize}
        \item Standard Level (SL): 35\%
        \item Higher Level (HL): 25\%
    \end{itemize}
\end{itemize}

\subsection*{Assessment Focus}
Paper 1 evaluates:
\begin{itemize}
    \item Understanding of core psychological approaches
    \item Ability to apply knowledge to contexts
    \item Conceptual thinking and integration
    \item Analytical and evaluative skills
\end{itemize}

\subsection*{Core Approaches}
\begin{itemize}
    \item Biological
    \item Cognitive
    \item Sociocultural
\end{itemize}

\subsection*{Key Concepts}
Students are expected to engage with:
\begin{itemize}
    \item Bias
    \item Causality
    \item Change
    \item Measurement
    \item Perspective
    \item Responsibility
\end{itemize}

\subsection*{Contexts}
Questions may be situated within:
\begin{itemize}
    \item Health and well-being
    \item Human development
    \item Human relationships
    \item Learning and cognition
\end{itemize}

\newpage

% ------------------------------------------------
\section{Section A: Short-Answer Questions}

\textbf{Marks Allocation:} 4 + 4

\subsection*{Structure}
\begin{itemize}
    \item Two compulsory questions
    \item Each drawn from different psychological approaches
    \item Focus on explanation of theory or concept
\end{itemize}

\subsection*{Skills Assessed}
\begin{itemize}
    \item Knowledge and understanding (AO1)
    \item Basic application (AO2)
\end{itemize}

\subsection*{High-Quality Response Characteristics}
\begin{itemize}
    \item Clear definition of key terms
    \item Accurate explanation of theory or concept
    \item Relevant and explained example
    \item Precise psychological terminology
\end{itemize}

\subsection*{Performance Bands (Generalized)}

\begin{longtable}{p{1.5cm} p{12cm}}
\toprule
\textbf{Marks} & \textbf{Performance Description} \\
\midrule
0 & Does not meet minimum requirements. \\
1--2 & Limited explanation; example may lack development. \\
3--4 & Clear and accurate explanation; relevant and well-linked example. \\
\bottomrule
\end{longtable}

\newpage

% ------------------------------------------------
\section{Section B: Application Questions}

\textbf{Marks Allocation:} 6 + 6

\subsection*{Structure}
\begin{itemize}
    \item Two compulsory scenario-based questions
    \item Each linked to a different context
    \item Requires application of theory to unseen material
\end{itemize}

\subsection*{Skills Assessed}
\begin{itemize}
    \item Application and analysis (AO2)
\end{itemize}

\subsection*{High-Quality Response Characteristics}
\begin{itemize}
    \item Accurate explanation of relevant theory
    \item Direct reference to scenario elements
    \item Logical and structured reasoning
    \item Clear connection between theory and context
\end{itemize}

\subsection*{Performance Bands (Generalized)}

\begin{longtable}{p{1.5cm} p{12cm}}
\toprule
\textbf{Marks} & \textbf{Performance Description} \\
\midrule
0 & Does not meet minimum requirements. \\
1--2 & Limited knowledge and weak application. \\
3--4 & Mostly accurate knowledge; partially developed application. \\
5--6 & Accurate and detailed knowledge; well-developed and explicit application. \\
\bottomrule
\end{longtable}

\newpage

% ------------------------------------------------
\section{Section C: Extended Response Question}

\textbf{Marks Allocation:} 15

\subsection*{Structure}
\begin{itemize}
    \item Two concept-based questions
    \item Students answer one
    \item Each linked to a specific context
\end{itemize}

\subsection*{Skills Assessed}
\begin{itemize}
    \item Knowledge and understanding (AO1)
    \item Synthesis and evaluation (AO3)
\end{itemize}

\subsection*{Expectations for High-Level Essays}
\begin{itemize}
    \item Sustained focus on the command term (e.g., evaluate, discuss)
    \item Integrated use of at least one key concept
    \item Development of a structured argument
    \item Critical analysis (strengths, limitations, implications)
    \item Clear and reasoned conclusion
\end{itemize}

\subsection*{Performance Bands (Generalized)}

\begin{longtable}{p{1.5cm} p{12cm}}
\toprule
\textbf{Marks} & \textbf{Performance Description} \\
\midrule
1--3 & Minimal understanding; largely descriptive. \\
4--6 & Some understanding; limited analytical depth. \\
7--9 & Adequate explanation; developing analysis. \\
10--12 & Clear explanation; relevant and consistent analysis. \\
13--15 & Comprehensive explanation; well-developed critical analysis; precise terminology and coherent conclusion. \\
\bottomrule
\end{longtable}

\newpage

% ------------------------------------------------
\section{Pedagogical Implementation Strategy}

\subsection*{Section A Preparation}
\begin{itemize}
    \item Focused theory drills
    \item Concept-definition practice
    \item Timed short responses
\end{itemize}

\subsection*{Section B Preparation}
\begin{itemize}
    \item Scenario-based workshops
    \item Explicit application modelling
    \item Context rotation exercises
\end{itemize}

\subsection*{Section C Preparation}
\begin{itemize}
    \item Concept mapping
    \item Evaluation paragraph training
    \item Structured essay planning under timed conditions
\end{itemize}

\vfill

\begin{center}
\rule{\textwidth}{0.4pt} \\
\small{This document is an internally developed academic resource. 
It summarizes assessment structure for instructional purposes and is not affiliated with or endorsed by the International Baccalaureate Organization.}
\end{center}

\end{document}